\subsection{Modular $A_4$ invariance and neutrino mixing --- arXiv:1808.03012}

\textbf{Authors:} Tatsuo Kobayashi, Naoya Omoto, Yusuke Shimizu, Kenta Takagi, Morimitsu Tanimoto, Takuya H. Tatsuishi

\paragraph{主な主張}
flavonを導入しないモジュラー $A_4$ 模型で、seesaw・Weinberg演算子・Dirac型のニュートリノ質量生成を比較し、当時の振動データと整合する模型と混合・CP位相の相関を示した\cite{Kobayashi:2018scp}。

\paragraph{新規性}
単一モジュライで制御される質量行列の現象論を質量生成機構ごとに検証し、ニュートリノ階層に応じた識別可能な予言を与えた。

\paragraph{位置付け}
モジュラー $A_4$ 模型の基本構成をニュートリノ観測と照合した初期の系統的研究である。
